\section{Paper Summary: FAS-Based IWIP Tracking}

\begin{frame}{论文概览}
    \small
    \textbf{论文:} Tracking Control for Inertia Wheel Inverted Pendulum Based on Fully Actuated System Theory\\[0.4em]
    \textbf{作者:} Haowen Liu, Shenjia Chen, Ping Li, Guangren Duan\\[0.6em]
    \begin{itemize}
        \setlength{\itemsep}{0.45em}
        \item 研究对象: 动量轮倒立摆（IWIP）的轨迹跟踪控制
        \item 核心思路: 将欠驱动系统通过微分同胚变换改写成高阶全驱系统（FAS）
        \item 主要目标: 在非线性模型下实现比局部线性控制更稳定的跟踪性能
        \item 主要增强: 引入可行参考轨迹约束，并设计扰动观测器进行在线补偿
    \end{itemize}
\end{frame}

\begin{frame}{系统模型与误差定义}
    \small
    论文从 Euler-Lagrange 模型出发：
    \[
    D(q)\ddot q + g(q) = u + d,
    \quad
    q = \begin{bmatrix} q_1 \\ q_2 \end{bmatrix},
    \quad
    u = \begin{bmatrix} 0 \\ \tau \end{bmatrix}
    \]
    其中 $q_1$ 为摆杆角，$q_2$ 为飞轮角，$d=[d_1,d_2]^T$ 为外扰。\\[0.5em]
    论文采用两条假设：
    \begin{itemize}
        \setlength{\itemsep}{0.4em}
        \item $q_1 \in \left(-\frac{\pi}{2}, \frac{\pi}{2}\right)$，即摆杆始终位于上半平面
        \item 扰动变化缓慢，满足 $\dot d \approx 0$
    \end{itemize}
    定义跟踪误差
    \[
    e_1 = q_{d1} - q_1,
    \qquad
    e_2 = q_{d2} - q_2,
    \]
    用于后续构造 FAS 输出变量。
\end{frame}

\begin{frame}{FAS 变量构造}
    \scriptsize
    论文定义新的全驱输出变量
    \[
    y_1 = \mu e_1 + e_2,
    \qquad
    \mu = \frac{D_{11}}{D_{12}}.
    \]
    其中
    \[
    y_{d1} = \mu q_{d1} + q_{d2},
    \qquad
    c_1 = (1-\mu)M,
    \qquad
    c_2 = b_1 - b_2.
    \]
    在 $\dot d_1 = 0$ 的假设下，原论文中的导数链可写为
    \[
    \begin{aligned}
    y_1 &= y_{d1} - \bigl(\mu q_1+q_2\bigr), \\
    \dot y_1 &= \dot y_{d1} - \bigl(\mu \dot q_1+\dot q_2\bigr), \\
    \ddot y_1 &= \ddot y_{d1} + c_1 \sin(q_1) + c_2 d_1, \\
    y_1^{(3)} &= y_{d1}^{(3)} + c_1 \cos(q_1)\dot q_1, \\
    y_1^{(4)} &= y_{d1}^{(4)} + c_1\!\left[
    -\sin(q_1)\dot q_1^{\,2}
    + \cos(q_1)\ddot q_1 \right].
    \end{aligned}
    \]
    通过消去 $q_1,\dot q_1$，原欠驱动系统可改写为四阶 SISO FAS：
    \[
    y_1^{(4)} = f(\ddot y_1,\dot y_1) + b(\ddot y_1)u_1.
    \]
\end{frame}

\begin{frame}{带偏移的模型构造}
    \scriptsize
    为了和原论文结果做对比，额外考虑摆角存在常值偏移 $\Delta$：
    \[
    \ddot q_1 = M\sin(q_1+\Delta)-b_1\tau+b_1d_1-b_1d_2,
    \quad
    \ddot q_2 = -M\sin(q_1+\Delta)+b_2\tau-b_1d_1+b_2d_2.
    \]
    此时误差与 FAS 输出改写为
    \[
    e_1 = q_{d1}-(q_1+\Delta),
    \qquad
    e_2 = q_{d2}-q_2,
    \qquad
    y_1 = \mu e_1 + e_2.
    \]
    若 $\Delta$ 为常值，则
    \[
    \sin(q_1+\Delta)=\sin q_1\cos\Delta+\cos q_1\sin\Delta,
    \]
    因而原论文中所有由 $\sin q_1,\cos q_1$ 组成的重力项，都需要整体替换为带偏移角的形式。
\end{frame}

\begin{frame}{带偏移的 $y_1$ 导数链}
    \scriptsize
    利用 $\mu=b_2/b_1$ 消去输入项，并采用 $\dot d_1=0$、$\dot\Delta=0$，有
    \[
    \begin{aligned}
    y_1 &= y_{d1} - \bigl(\mu(q_1+\Delta)+q_2\bigr), \\
    \dot y_1 &= \dot y_{d1} - \bigl(\mu \dot q_1+\dot q_2\bigr), \\
    \ddot y_1 &= \ddot y_{d1} + c_1 \sin(q_1+\Delta) + c_2 d_1, \\
    y_1^{(3)} &= y_{d1}^{(3)} + c_1 \cos(q_1+\Delta)\dot q_1, \\
    y_1^{(4)} &= y_{d1}^{(4)} + c_1\!\left[
    -\sin(q_1+\Delta)\dot q_1^{\,2}
    + \cos(q_1+\Delta)\ddot q_1 \right].
    \end{aligned}
    \]
    通过消去 $q_1,\dot q_1$，原欠驱动系统可改写为四阶 SISO FAS：
    \[
    y_1^{(4)} = f(\ddot y_1,\dot y_1) + b(\ddot y_1)u_1.
    \]
\end{frame}

\begin{frame}{无偏移与有偏移的对比}
    \scriptsize
    \begin{columns}[T]
        \begin{column}{0.48\textwidth}
            \textbf{原论文}
            \[
            e_1 = q_{d1} - q_1
            \]
            \[
            \ddot y_1 = \ddot y_{d1} + c_1 \sin(q_1) + c_2 d_1
            \]
            \[
            y_1^{(3)} = y_{d1}^{(3)} + c_1 \cos(q_1)\dot q_1
            \]
            \[
            \begin{aligned}
            y_1^{(4)} &= y_{d1}^{(4)} \\
            &\quad + c_1\!\left[
            -\sin(q_1)\dot q_1^2
            + \cos(q_1)\ddot q_1
            \right]
            \end{aligned}
            \]
        \end{column}
        \begin{column}{0.48\textwidth}
            \textbf{考虑偏移 $\Delta$}
            \[
            e_1 = q_{d1} - (q_1+\Delta)
            \]
            \[
            \ddot y_1 = \ddot y_{d1} + c_1 \sin(q_1+\Delta) + c_2 d_1
            \]
            \[
            y_1^{(3)} = y_{d1}^{(3)} + c_1 \cos(q_1+\Delta)\dot q_1
            \]
            \[
            \begin{aligned}
            y_1^{(4)} &= y_{d1}^{(4)} \\
            &\quad + c_1\!\left[
            -\sin(q_1+\Delta)\dot q_1^2 \right.\\
            &\qquad\left. + \cos(q_1+\Delta)\ddot q_1
            \right]
            \end{aligned}
            \]
        \end{column}
    \end{columns}
    偏移不会改变 FAS 的设计框架，但会改变误差定义和所有重力非线性项。
\end{frame}

\begin{frame}{名义 FAS 控制律设计}
    \scriptsize
    原文先将 IWIP 写成四阶 FAS
    \[
    y_1^{(4)} = f(\ddot y_1,\dot y_1) + b(\ddot y_1)u_1.
    \]
    为了指定期望误差动力学，选取名义控制律
    \[
    u_1 =
    -\frac{1}{b(\ddot y_1)}
    \left[
    f(\ddot y_1,\dot y_1) + \sum_{i=0}^{3} K_i y_1^{(i)}
    \right].
    \]
    代回 FAS 可得闭环
    \[
    y_1^{(4)} + K_3 y_1^{(3)} + K_2 \ddot y_1 + K_1 \dot y_1 + K_0 y_1 = 0.
    \]
    因此特征多项式取为
    \[
    \lambda_c(s)=s^4+K_3 s^3+K_2 s^2+K_1 s+K_0.
    \]
    按原文采用极点配置
    \[
    \lambda_c(s)=(s+\omega_c)^4
    = s^4 + 4\omega_c s^3 + 6\omega_c^2 s^2 + 4\omega_c^3 s + \omega_c^4,
    \]
    从而得到
    \[
    K_0=\omega_c^4,\qquad
    K_1=4\omega_c^3,\qquad
    K_2=6\omega_c^2,\qquad
    K_3=4\omega_c.
    \]
\end{frame}

\begin{frame}{可行参考轨迹约束}
    \scriptsize
    原文指出该 FAS 实际上是 SISO 结构，因此 $q_{d1}(t),q_{d2}(t)$ 不能完全独立指定。\\[0.4em]
    当 $y_1 \to 0$ 且 $\ddot y_1 \to 0$ 时，有
    \[
    y_1 = \mu e_1 + e_2 = 0,
    \qquad
    \ddot y_1 = \ddot y_{d1} + c_1 \sin(q_1) + c_2 d_1 = 0.
    \]
    为保证零误差跟踪，期望轨迹必须满足
    \[
    \ddot y_{d1} = -c_1 \sin(q_{d1}) - c_2 d_1.
    \]
    代回 $y_{d1}=\mu q_{d1}+q_{d2}$ 并利用参数定义，可写成论文中的约束式
    \[
    D_{11}\ddot q_{d1}+D_{12}\ddot q_{d2}-\bar m g \sin(q_{d1})=d_1.
    \]
    这个条件的物理含义是：
    \begin{itemize}
        \setlength{\itemsep}{0.3em}
        \item 期望摆角和飞轮轨迹必须与系统动力学一致
        \item 参考输入的自由度被约束到一维可实现流形上
        \item 若不满足该约束，控制器只能做到有界误差而非严格零误差
    \end{itemize}
\end{frame}

\begin{frame}{带扰动补偿的实际控制律}
    \scriptsize
    将扰动估计并入控制后，原文给出的实际输入力矩为
    \[
    \tau=
    \frac{1}{b_1 b(\ddot y_1)}
    \left[
    f(\ddot y_1,\dot y_1) + \sum_{i=0}^{3} K_i \hat y_1^{(i)}
    \right]
    +\frac{M}{b_1}\sin(q_1)+\hat d_1-\hat d_2.
    \]
    代回子 FAS 后得到受扰闭环
    \[
    y_1^{(4)}+\sum_{i=0}^{3}K_i y_1^{(i)}+\Delta(t,y_1,\dot y_1,\ddot y_1)=0,
    \]
    其中扰动项为
    \[
    \Delta=(b_1\tilde d_1-b_1\tilde d_2)b(\ddot y_1)-K_2\tilde d_1.
    \]
    实际实现时，$\hat d$ 由扰动观测器在线给出，用于补偿约束失配和外部扰动。
\end{frame}

\begin{frame}{考虑偏移 $\Delta$ 的名义控制律}
    \scriptsize
    若前面建立的 FAS 变量改为
    \[
    e_1=q_{d1}-(q_1+\Delta),
    \qquad
    y_1=\mu e_1+e_2,
    \]
    且 $\Delta$ 为常值，则导数链中的重力项全部替换为
    \[
    \sin(q_1+\Delta),\qquad \cos(q_1+\Delta).
    \]
    此时四阶 FAS 仍写成同一结构
    \[
    y_1^{(4)} = f_\Delta(\ddot y_1,\dot y_1) + b_\Delta(\ddot y_1)u_{1,\Delta}.
    \]
    因而名义控制律可以平行写为
    \[
    u_{1,\Delta}
    =
    -\frac{1}{b_\Delta(\ddot y_1)}
    \left[
    f_\Delta(\ddot y_1,\dot y_1)+\sum_{i=0}^{3}K_i y_1^{(i)}
    \right].
    \]
    对应闭环仍然是
    \[
    y_1^{(4)} + K_3 y_1^{(3)} + K_2 \ddot y_1 + K_1 \dot y_1 + K_0 y_1 = 0,
    \]
    所以极点配置不需要改变，仍取
    \[
    \lambda_c(s)=(s+\omega_c)^4,
    \quad
    K_0=\omega_c^4,\;
    K_1=4\omega_c^3,\;
    K_2=6\omega_c^2,\;
    K_3=4\omega_c.
    \]
\end{frame}

\begin{frame}{考虑偏移 $\Delta$ 的实际力矩控制律}
    \scriptsize
    由
    \[
    \ddot y_1 = \ddot y_{d1} + c_1\sin(q_1+\Delta)+c_2 d_1
    \]
    可见在执行层，原文中的补偿项
    \[
    \frac{M}{b_1}\sin(q_1)
    \]
    应同步改为
    \[
    \frac{M}{b_1}\sin(q_1+\Delta).
    \]
    因此带扰动观测器的实际输入力矩可写成
    \[
    \tau_\Delta=
    \frac{1}{b_1 b_\Delta(\ddot y_1)}
    \left[
    f_\Delta(\ddot y_1,\dot y_1)+\sum_{i=0}^{3}K_i \hat y_1^{(i)}
    \right]
    +\frac{M}{b_1}\sin(q_1+\Delta)+\hat d_1-\hat d_2.
    \]
    若进一步展开偏移角，则
    \[
    \tau_\Delta=
    \frac{1}{b_1 b_\Delta(\ddot y_1)}
    \left[
    f_\Delta+\sum_{i=0}^{3}K_i \hat y_1^{(i)}
    \right]
    +\frac{M}{b_1}\bigl(\sin q_1\cos\Delta+\cos q_1\sin\Delta\bigr)
    +\hat d_1-\hat d_2.
    \]
    这说明偏移不会改变控制器设计框架，只会改变：
    \begin{itemize}
        \setlength{\itemsep}{0.3em}
        \item FAS 映射中的非线性函数 $f,b$
        \item 力矩控制中的重力补偿项
        \item 可行参考轨迹约束中的正弦项
    \end{itemize}
\end{frame}

\begin{frame}{论文结论与可借鉴点}
    \small
    \begin{itemize}
        \setlength{\itemsep}{0.5em}
        \item 论文的核心贡献不是重新推导 IWIP 模型，而是把 IWIP 提升为可直接设计的四阶 FAS
        \item 参考轨迹约束揭示了欠驱动系统只能跟踪\textbf{可实现}的目标，而非两个完全独立的输出
        \item 扰动观测器补偿后，仿真与实验都表明跟踪性能和鲁棒性得到改善
        \item 对你当前这份 beamer 而言，这篇论文最值得保留的主线是：
        “原系统建模 $\rightarrow$ FAS 变换 $\rightarrow$ 四阶控制器 $\rightarrow$ 轨迹约束 $\rightarrow$ 扰动观测器”
    \end{itemize}
\end{frame}
